\chapter{Tricyclic Antidepressant (TCA) Screen}
\section*{What Is This Test?} A TCA screen detects tricyclic antidepressants or related compounds in a specimen, usually during suspected overdose evaluation.\section*{Alternative Names} TCA toxicology screen; tricyclic drug level; antidepressant overdose screen.\section*{Principle} Initial immunoassays may cross-react and are generally qualitative; targeted chromatography can identify specific drugs.\section*{Why This Test Is Ordered} It supports assessment of unexplained coma, seizures, arrhythmia, or suspected TCA ingestion.\section*{Primary vs Secondary Test} ECG, airway, circulation, and clinical toxicology assessment are primary; the screen does not rule out poisoning.\section*{How This Test Is Performed} Blood or urine is collected and tested according to the laboratory method; ECG and electrolytes are obtained urgently when overdose is possible.\section*{What Preparation Is Needed} None in an emergency. Provide medicine, dose, and timing information if known.\section*{Understanding Results} A positive screen may reflect prescribed medicine or cross-reactivity and does not measure severity; a negative result does not exclude toxicity.\section*{Follow-up Tests} ECG, electrolytes, blood gas, glucose, renal function, confirmatory toxicology, and poison-center consultation may follow.\section*{References} \begin{enumerate}\item MedlinePlus. Tricyclic Antidepressant (TCA) Screen.\item American College of Medical Toxicology. TCA poisoning guidance.\end{enumerate}\clearpage\section*{Understanding Results and Follow-up}
The laboratory report should identify the specimen, method, units, reference interval or decision limit, and whether the result is positive, negative, abnormal, or inconclusive. Reference intervals vary with age, sex, pregnancy, specimen type, assay, and laboratory; a result outside a range is not automatically a diagnosis, and a result within a range does not always exclude disease. Discuss unexpected results with the ordering clinician, who can decide whether repeat or confirmatory testing, treatment, or referral is appropriate. Do not change medicines or delay urgent care based on an isolated result.
\section*{Regulatory Status and Authorization Date}This chapter describes a test category, not a specific commercial product. FDA, CDSCO, EU IVDR, and MHRA approval or registration dates therefore cannot be assigned without the assay manufacturer, product name, instrument, and intended use. For laboratory-developed tests, an individual agency approval date may not exist. See the Preface for the interpretation of regulatory status.
\clearpage
